% Options for packages loaded elsewhere
\PassOptionsToPackage{unicode}{hyperref}
\PassOptionsToPackage{hyphens}{url}
\documentclass[
]{article}
\usepackage{xcolor}
\usepackage{amsmath,amssymb}
\setcounter{secnumdepth}{-\maxdimen} % remove section numbering
\usepackage{iftex}
\ifPDFTeX
  \usepackage[T1]{fontenc}
  \usepackage[utf8]{inputenc}
  \usepackage{textcomp} % provide euro and other symbols
\else % if luatex or xetex
  \usepackage{unicode-math} % this also loads fontspec
  \defaultfontfeatures{Scale=MatchLowercase}
  \defaultfontfeatures[\rmfamily]{Ligatures=TeX,Scale=1}
\fi
\usepackage{lmodern}
\ifPDFTeX\else
  % xetex/luatex font selection
\fi
% Use upquote if available, for straight quotes in verbatim environments
\IfFileExists{upquote.sty}{\usepackage{upquote}}{}
\IfFileExists{microtype.sty}{% use microtype if available
  \usepackage[]{microtype}
  \UseMicrotypeSet[protrusion]{basicmath} % disable protrusion for tt fonts
}{}
\makeatletter
\@ifundefined{KOMAClassName}{% if non-KOMA class
  \IfFileExists{parskip.sty}{%
    \usepackage{parskip}
  }{% else
    \setlength{\parindent}{0pt}
    \setlength{\parskip}{6pt plus 2pt minus 1pt}}
}{% if KOMA class
  \KOMAoptions{parskip=half}}
\makeatother
\usepackage{color}
\usepackage{fancyvrb}
\newcommand{\VerbBar}{|}
\newcommand{\VERB}{\Verb[commandchars=\\\{\}]}
\DefineVerbatimEnvironment{Highlighting}{Verbatim}{commandchars=\\\{\}}
% Add ',fontsize=\small' for more characters per line
\newenvironment{Shaded}{}{}
\newcommand{\AlertTok}[1]{\textcolor[rgb]{1.00,0.00,0.00}{\textbf{#1}}}
\newcommand{\AnnotationTok}[1]{\textcolor[rgb]{0.38,0.63,0.69}{\textbf{\textit{#1}}}}
\newcommand{\AttributeTok}[1]{\textcolor[rgb]{0.49,0.56,0.16}{#1}}
\newcommand{\BaseNTok}[1]{\textcolor[rgb]{0.25,0.63,0.44}{#1}}
\newcommand{\BuiltInTok}[1]{\textcolor[rgb]{0.00,0.50,0.00}{#1}}
\newcommand{\CharTok}[1]{\textcolor[rgb]{0.25,0.44,0.63}{#1}}
\newcommand{\CommentTok}[1]{\textcolor[rgb]{0.38,0.63,0.69}{\textit{#1}}}
\newcommand{\CommentVarTok}[1]{\textcolor[rgb]{0.38,0.63,0.69}{\textbf{\textit{#1}}}}
\newcommand{\ConstantTok}[1]{\textcolor[rgb]{0.53,0.00,0.00}{#1}}
\newcommand{\ControlFlowTok}[1]{\textcolor[rgb]{0.00,0.44,0.13}{\textbf{#1}}}
\newcommand{\DataTypeTok}[1]{\textcolor[rgb]{0.56,0.13,0.00}{#1}}
\newcommand{\DecValTok}[1]{\textcolor[rgb]{0.25,0.63,0.44}{#1}}
\newcommand{\DocumentationTok}[1]{\textcolor[rgb]{0.73,0.13,0.13}{\textit{#1}}}
\newcommand{\ErrorTok}[1]{\textcolor[rgb]{1.00,0.00,0.00}{\textbf{#1}}}
\newcommand{\ExtensionTok}[1]{#1}
\newcommand{\FloatTok}[1]{\textcolor[rgb]{0.25,0.63,0.44}{#1}}
\newcommand{\FunctionTok}[1]{\textcolor[rgb]{0.02,0.16,0.49}{#1}}
\newcommand{\ImportTok}[1]{\textcolor[rgb]{0.00,0.50,0.00}{\textbf{#1}}}
\newcommand{\InformationTok}[1]{\textcolor[rgb]{0.38,0.63,0.69}{\textbf{\textit{#1}}}}
\newcommand{\KeywordTok}[1]{\textcolor[rgb]{0.00,0.44,0.13}{\textbf{#1}}}
\newcommand{\NormalTok}[1]{#1}
\newcommand{\OperatorTok}[1]{\textcolor[rgb]{0.40,0.40,0.40}{#1}}
\newcommand{\OtherTok}[1]{\textcolor[rgb]{0.00,0.44,0.13}{#1}}
\newcommand{\PreprocessorTok}[1]{\textcolor[rgb]{0.74,0.48,0.00}{#1}}
\newcommand{\RegionMarkerTok}[1]{#1}
\newcommand{\SpecialCharTok}[1]{\textcolor[rgb]{0.25,0.44,0.63}{#1}}
\newcommand{\SpecialStringTok}[1]{\textcolor[rgb]{0.73,0.40,0.53}{#1}}
\newcommand{\StringTok}[1]{\textcolor[rgb]{0.25,0.44,0.63}{#1}}
\newcommand{\VariableTok}[1]{\textcolor[rgb]{0.10,0.09,0.49}{#1}}
\newcommand{\VerbatimStringTok}[1]{\textcolor[rgb]{0.25,0.44,0.63}{#1}}
\newcommand{\WarningTok}[1]{\textcolor[rgb]{0.38,0.63,0.69}{\textbf{\textit{#1}}}}
\setlength{\emergencystretch}{3em} % prevent overfull lines
\providecommand{\tightlist}{%
  \setlength{\itemsep}{0pt}\setlength{\parskip}{0pt}}
\usepackage{bookmark}
\IfFileExists{xurl.sty}{\usepackage{xurl}}{} % add URL line breaks if available
\urlstyle{same}
% fallback for those not using the hyperref driver hyperxmp:
\makeatletter
\@ifundefined{xmpquote}{\newcommand{\xmpquote}[1]{#1}}{}
\makeatother
\hypersetup{
  hidelinks,
  pdfcreator={LaTeX via pandoc}}

\author{}
\date{}

\begin{document}

{
\setcounter{tocdepth}{2}
\tableofcontents
}
\section{WebCoder --- Contribuții}\label{webcoder--contribuux21bii}

\begin{quote}
Rezumat într-un paragraf: WebCoder este un monorepo pe o singură ramură
(\texttt{main}) --- Next.js 16 la rădăcină, Django 6 sub
\texttt{backend/} --- în care orice modificare ajunge printr-un PR care
ține VERZI ambele baterii de teste (frontend: vitest + tsc + lint fără
warninguri + build + Playwright; backend: pytest cu 100 \% acoperire de
linii și fără warninguri). Ghidul explică porțile exacte, convențiile de
PR, fluxul BACKLOG și cum au funcționat „valurile" de roadmap.
\end{quote}

\textbf{Cui se adresează:} oricui contribuie cod, conținut sau
documentație. Perechea în engleză:
\href{../guide-en/CONTRIBUTING.md}{\texttt{docs/guide-en/CONTRIBUTING.md}}.

\begin{center}\rule{0.5\linewidth}{0.5pt}\end{center}

\subsection{1. Regulile
repository-ului}\label{1-regulile-repository-ului}

\begin{itemize}
\tightlist
\item
  \textbf{\texttt{main} este singura ramură de bază.} Vechiul split
  (ramura \texttt{backend}) a fost unificat în octombrie 2026; fiecare
  PR țintește \texttt{main}. Nu redeschideți splitul.
\item
  Repo-ul e monorepo: frontend la rădăcină (\texttt{src/},
  \texttt{package.json}), backend sub \texttt{backend/}, e2e sub
  \texttt{e2e/}, shell-ul mobil sub \texttt{mobile/}.
\item
  Merge-ul implicit este \textbf{squash}
  (\texttt{gh\ pr\ merge\ -\/-squash\ -\/-delete-branch}); commit-urile
  de merge adevărate sunt rezervate operațiunilor care păstrează
  istoric.
\item
  Lucrul se face în worktree-uri de unică folosință --- niciodată nu se
  comite într-un checkout de deploy live.
\end{itemize}

\subsection{2. Porțile de calitate (obligatorii înainte de
push)}\label{2-porux21bile-de-calitate-obligatorii-uxeenainte-de-push}

Frontend --- din rădăcina repo-ului:

\begin{Shaded}
\begin{Highlighting}[]
\ExtensionTok{npm}\NormalTok{ run test                      }\CommentTok{\# vitest}
\ExtensionTok{npx}\NormalTok{ tsc }\AttributeTok{{-}p}\NormalTok{ tsconfig.json }\AttributeTok{{-}{-}noEmit} \CommentTok{\# typecheck}
\ExtensionTok{npm}\NormalTok{ run lint                      }\CommentTok{\# 0 erori ȘI 0 warninguri (politica zero{-}warninguri)}
\ExtensionTok{npm}\NormalTok{ run build                     }\CommentTok{\# next build}
\ExtensionTok{npx}\NormalTok{ playwright test               }\CommentTok{\# e2e, mod mock{-}api (ermetic)}
\end{Highlighting}
\end{Shaded}

Backend --- orice interpretor cu \texttt{backend/deps.txt} instalat:

\begin{Shaded}
\begin{Highlighting}[]
\VariableTok{SECRET\_KEY}\OperatorTok{=}\NormalTok{x }\VariableTok{DB\_PASSWORD}\OperatorTok{=}\NormalTok{x }\ExtensionTok{python} \AttributeTok{{-}m}\NormalTok{ pytest }\AttributeTok{{-}q}    \CommentTok{\# profil CI sqlite, 100 \% acoperire, 0 warninguri}
\end{Highlighting}
\end{Shaded}

Reguli de calitate aplicate peste:

\begin{itemize}
\tightlist
\item
  \textbf{Orice modul backend NOU atinge 100 \% acoperire de linii} ---
  bara repo-wide a fost 100 \% înainte; mențineți-o.
\item
  \textbf{Zero warninguri}: suita pytest rulează curat (rezolvați la
  sursă sau adăugați un ignore \texttt{filterwarnings} îngust, comentat
  individual); \texttt{npm\ run\ lint} raportează zero probleme.
\item
  Ruff format + biome format rulează în CI; codul neformatat cade la
  poarta de calitate.
\item
  Migrațiile nu trebuie să se colaboreze: valurile numerotează
  migrațiile independent --- dacă două ramuri adaugă migrații cu același
  număr, renumerotați-o pe cea mai tânără înainte de merge (vezi
  PR-urile \#360, \#375).
\end{itemize}

\subsection{3. Convențiile de PR}\label{3-convenux21biile-de-pr}

\begin{enumerate}
\def\labelenumi{\arabic{enumi}.}
\tightlist
\item
  Ramifică din \texttt{origin/main}
  (\texttt{git\ worktree\ add\ \textless{}cale\textgreater{}\ -b\ \textless{}ramura\textgreater{}\ origin/main}).
\item
  Implementează \textbf{cu teste} --- pytest backend pentru fiecare
  endpoint/comandă nouă; vitest pentru componente; acoperire Playwright
  pentru fluxurile vizibile pentru utilizator (extindeți fixture-urile
  \texttt{e2e/mock-api}; adăugați specificații real-stack când
  funcționalitatea atinge evaluatorul sau autentificarea).
\item
  Verifică local (porțile exacte de mai sus), push, deschide PR-ul către
  \texttt{main} cu \texttt{gh\ pr\ create}.
\item
  Urmărește check-urile până la verde (CI rulează \texttt{frontend.yml}
  pe căile din rădăcină, \texttt{django.yml} pe \texttt{backend/**},
  CodeQL pe tot) --- până la 20 × 30 s.
\item
  \texttt{gh\ pr\ merge\ -\/-squash\ -\/-delete-branch}, apoi șterge
  worktree-ul.
\item
  \textbf{Validare pe stack-ul live după merge-uri de backend} (RUNBOOK
  §convenții PR):
  \texttt{git\ -C\ \textless{}worktree\ live\textgreater{}\ pull\ -\/-ff-only\ origin\ main},
  rebuild \texttt{backend\ celery-worker\ judge-runner},
  \texttt{manage.py\ migrate}.
\end{enumerate}

Stil de commit: \texttt{tip(scop):\ rezumat} --- ex.
\texttt{feat(contests):\ …}, \texttt{fix(e2e):\ …},
\texttt{docs(guides):\ …}. PR-urile valurilor sunt o referință bună
(\texttt{git\ log\ -\/-oneline}).

\subsection{4. Fluxul BACKLOG}\label{4-fluxul-backlog}

\texttt{docs/BACKLOG.md} este tracker-ul de nivel element, de referință.
Reguli (stated la începutul fișierului):

\begin{itemize}
\tightlist
\item
  statusuri: \texttt{shipped\ (PR\ \#N)} · \texttt{in-flight\ (wave\ X)}
  · \texttt{planned\ (wave\ E/F/G)} · \texttt{someday};
\item
  \textbf{agentul valului care livrează un element îi actualizează
  statusul în același PR};
\item
  numerotarea e stabilă --- nu renumerotați; retrageți cu
  \texttt{\textasciitilde{}}.
\end{itemize}

Deasupra stau \texttt{docs/SPEC-GAPS.md} (backlog-ul de platformă cu
id-urile B/J/F/A/Q/SP/X) și
\texttt{docs/ROADMAP-COMPETITIVE-2026-10-22.md} (catalogul de 78 gap-uri
+ secvențierea valurilor + corecțiunile binding din round-2). Ideile noi
merg în parking lot-ul BACKLOG; gap-urile noi de platformă primesc un
rând în SPEC-GAPS.

\subsection{5. Cum au funcționat
valurile}\label{5-cum-au-funcux21bionat-valurile}

Programul de feature-uri 2026-11 a rulat ca valuri secvențiale (0 → A →
B → C → D → E → F → G), fiecare un set de PR-uri care implementau o
felie de roadmap:

\begin{enumerate}
\def\labelenumi{\arabic{enumi}.}
\tightlist
\item
  Specificația valului se citea din ROADMAP §4 + adaosurile round-2
  (ops-înainte-de-feature-uri, moderare-înainte-de-UGC, porți legale ---
  acestea înving roadmap-ul când nu cad de accord pe ordine).
\item
  Elementele livrau ca grupuri coerente de PR-uri (ramură → teste →
  verificare locală → PR → check-uri → squash-merge), fiecare PR
  actualizând statusurile din BACKLOG.
\item
  Elementele care atingeau evaluatorul/autentificarea se validau
  suplimentar pe stack-ul compose live (RUNBOOK Mode A), inclusiv
  Playwright în mod real
  (\texttt{E2E\_API\_BASE=http://127.0.0.1:8010}).
\item
  Lucrul de conținut (importul de \textasciitilde1 850 de probleme,
  generarea testelor ascunse) a rulat ca management commands comise în
  repo (\texttt{catalog\_polish}, \texttt{boost\_hidden\_tests},
  \texttt{export\_catalog}/\texttt{import\_catalog}) --- astfel încât
  orice mutație de date e review-abilă și reproductibilă.
\end{enumerate}

\subsection{6. Contribuții de
conținut}\label{6-contribuux21bii-de-conux21binut}

\begin{itemize}
\tightlist
\item
  Problemele se author-esc \textbf{bilingv}: enunțul în engleză e
  obligatoriu (fallback universal), româna e puternic recomandată; încă
  10 limbi sunt opționale per problemă.
\item
  Sarcinile de tip Bac „scrieți subprogramul" se rescriu la import ca
  sarcini cu program complet I/O (aceeași exercițiu, cuvintele noastre
  --- vezi \texttt{problems/data/hidden\_test\_specs.py}).
\item
  Suite de teste: nu importați date de test third-party fără drepturi;
  generați suite ascunse prin formatul de specificații comis (soluție de
  referință + inputuri + verificare de flip pe soluție greșită).
\item
  Ledger-ul de permisiuni CNT-2 (\texttt{ImportSource}) înregistrează
  sursa, termenii și dovada pentru fiecare corpus importat.
\end{itemize}

\subsection{7. Contribuții de
documentație}\label{7-contribuux21bii-de-documentaux21bie}

Suita de docs e bilingvă (EN în \texttt{docs/guide-en/}, RO în
\texttt{docs/guide-ro/}) și se build-uiește în PDF/DOCX/TEX prin
\texttt{docs/build-docs.ps1} (pandoc + tectonic; outputurile se comit
sub \texttt{docs/dist/}). Când schimbați un comportament documentat
într-un ghid, actualizați ambele perechi lingvistice în același PR și
rebuild-uiți artefactele:

\begin{Shaded}
\begin{Highlighting}[]
\NormalTok{docs}\OperatorTok{/}\NormalTok{build{-}docs}\OperatorTok{.}\FunctionTok{ps1}    \CommentTok{\# regenerează docs/dist din sursele .md}
\end{Highlighting}
\end{Shaded}

Diagrame: sursele Mermaid stau în \texttt{docs/diagrams/*.mmd} cu
PNG-uri randate lângă ele --- editați \texttt{.mmd}, re-randați
(mermaid-cli), comiteți ambele.

\end{document}
